\subsection{Compact perturbation and essential spectrum}

In this number, E is an infinite-dimensional complex Hilbert space.

Lemma 13. --- Let I be an orthonormal family in E. The family I converges weakly to 0 along the filter of complements of finite subsets of I.

Let \(x \in E\) . By Bessel's inequality (EVT, V, p. 21, prop. 4) and TG, IV, p. 37, th. 1, the family \((|\langle e_{i} | x\rangle|^{2})_{i\in I}\) is summable, hence \(\langle e_{i} | x\rangle\) tends to 0 along the filter of complements of finite subsets of I (TG, III, p. 38, prop. 1).

PROPOSITION 18. --- Let \(u\) be a self-adjoint partial operator on E and let \(\lambda\) be a real number. We have \(\lambda \in \mathrm{Sp}_{\mathrm{e}}(u)\) if and only if there exists an orthonormal sequence \((x_{n})_{n \in \mathbf{N}}\) in E such that \(u(x_{n}) - \lambda x_{n}\) tends to 0 in E.

Suppose first that \(\lambda \in \mathrm{Sp}_{\mathrm{e}}(u)\) . If the spectral projector of \(u\) relative to \(\{\lambda\}\) is of infinite rank, every orthonormal sequence \((x_{n})\) in its image answers the question (cf. cor. of prop. 9 of IV, p. 278). In what follows we shall assume that this is not the case.

For every \(k \in \mathbf{N}\) , let \(\mathrm{J}_k\) be the set of \(t \in [\lambda - 1, \lambda + 1]\) such that

\[
{\frac {1}{k + 2}} <   | t - \lambda | \leqslant {\frac {1}{k + 1}}.
\]

The sets \(J_{k}\) are pairwise disjoint. Moreover, for every integer \(K \in N\) , the sets \((\mathrm{J}_{k})_{k \geqslant K}\) form a partition of the set

\[
\mathrm{I} _ {\mathrm{K}} = [ \lambda - 1 / (\mathrm{K} + 1), \lambda + 1 / (\mathrm{K} + 1) ] - \{0 \}.
\]

Since the spectral projector of \(u\) relative to \(\mathrm{I_K} \cup \{0\}\) has infinite rank (lemma 9 of IV, p. 297) and the one relative to \(\{\lambda\}\) is assumed to have finite rank, it follows from prop. 8 of IV, p. 278 that there exists a sequence \((k_n)_{n \in \mathbf{N}}\) strictly increasing in \(\mathbf{N}\) such that the spectral projector \(p_{n}\) of u relative to \(J_{k_{n}}\) is nonzero. Let \(x_{n}\) be a vector of norm 1 in the image of \(p_{n}\) . The sequence \((x_{n})\) is orthonormal, since the image of \(p_{n}\) is orthogonal to that of \(p_{m}\) for all \(n \neq m\) in N.

Let \(n \in N\) . Denote by \(\mu_{n}\) the spectral measure of \(x_{n}\) relative to u; its support is contained in \(J_{k_{n}}\) (prop. 9 of IV, p. 278). It follows

\[
\| u (x _ {n}) - \lambda x _ {n} \| ^ {2} = \int_ {\mathbf {C}} | t - \lambda | ^ {2} d \mu_ {n} (t) \leqslant \frac {1}{k _ {n} ^ {2}},
\]

therefore the sequence \((x_{n})_{n\in\mathbf{N}}\) has the required property.

Suppose conversely that there exists an orthonormal sequence \((x_{n})_{n\in \mathbf{N}}\) in E such that \(u(x_{n}) - \lambda x_{n}\to 0\) . Denote by \(\mu_{n}\) the spectral measure of \(x_{n}\) relative to \(u\) .

Let \(\varepsilon > 0\) . Denote by \(p_{\varepsilon}\) be the spectral projector of \(u\) relative to the open interval \(\mathrm{I}_{\varepsilon} = ]\lambda - \varepsilon, \lambda + \varepsilon[\) . For every \(n \in \mathbf{N}\) , it follows that

\[
\begin{array}{r l} & 1 = \| x _ {n} \| ^ {2} = \mu_ {n} (\mathrm{I} _ {\varepsilon}) + \mu_ {n} (\mathbf {R} - \mathrm{I} _ {\varepsilon}) \\ & \qquad \leqslant \mu_ {n} (\mathrm{I} _ {\varepsilon}) + \frac {1}{\varepsilon^ {2}} \int_ {\mathbf {R} - \mathrm{I} _ {\varepsilon}} | t - \lambda | ^ {2} d \mu_ {n} (t) \\ & \qquad = \| p _ {\varepsilon} (x _ {n}) \| ^ {2} + \frac {1}{\varepsilon^ {2}} \| u (x _ {n}) - \lambda x _ {n} \| ^ {2}. \end{array}
\]

The hypothesis on the sequence \((x_{n})\) therefore implies that the norm of \(p_{\varepsilon}(x_{n})\) cannot tend to 0. Since the orthonormal sequence \((x_{n})\) converges weakly to 0 in E (Lemma 13), the projector \(p_{\varepsilon}\) cannot be compact (cor. of prop. 6 of III, p. 6) and is consequently of infinite rank. Since this holds for every \(\varepsilon > 0\) , lemma 9 of IV, p. 297 allows us to conclude that \(\lambda \in \mathrm{Sp}_{\mathrm{e}}(u)\) .

The following theorem is the analogue for self-adjoint partial operators of Theorem 3 in III, p. 93.

THEOREM 4. --- Let u be a self-adjoint partial operator on E. The essential spectrum of u is the intersection of the sets Sp(u + v) where v ranges over the set of compact Hermitian endomorphisms of E.

The partial operator \(u + v\) is self-adjoint for every Hermitian endomorphism v of E since \((u + v)^{*} = u^{*} + v^{*}\) (cf. IV, p. 236).

Let us prove that if \(v\) is compact, then \(\mathrm{Sp_e}(u) \subset \mathrm{Sp_e}(u + v)\) . Let \(\lambda \in \mathrm{Sp_e}(u)\) and let \((x_n)_{n \in \mathbf{N}}\) be an orthonormal sequence in E such that \(u(x_n) - \lambda x_n\) converges to 0 (prop. 18). The sequence \((x_n)\) converges weakly to 0 in E (lemma 13), and since the endomorphism \(v\) is compact, the sequence \((v(x_n))_{n \in \mathbf{N}}\) converges to 0 in E (cor. of prop. 6 of III, p. 6). Consequently, the sequence \((u+v)(x_{n})-\lambda x_{n}\) converges to 0 in E, and Prop. 18 implies that \(\lambda\in\mathrm{Sp}_{\mathrm{e}}(u+v)\) .

The essential spectrum of \(u\) is therefore contained in the intersection of the sets \(\mathrm{Sp}(u + v)\) for Hermitian \(v \in \mathcal{L}^{\mathrm{c}}(\mathrm{E})\).

Conversely, let \(\lambda \in \mathrm{Sp}_{\mathrm{s}}(u)\) . Denote by \(\mathrm{E}_{\lambda}\) the eigenspace of \(u\) relative to \(\lambda\) , and \(p_{\lambda}\) the orthogonal projector with image \(\mathrm{E}_{\lambda}\) ; it is the spectral projector of \(u\) relative to \(\{\lambda\}\) (corollary of Prop. 9 of IV, p. 278). By definition of the sensitive spectrum, the projector \(p_{\lambda}\) is of finite rank, hence compact. Set \(w = u + p_{\lambda}\) ; it is a self-adjoint partial operator. To conclude the proof, let us verify that \(\lambda\) belongs to the resolvent set of \(w\) .

One has \(\mathrm{E}_{\lambda} \subset \operatorname{dom}(u)\) and the spaces \(\mathrm{E}_{\lambda}\) and \(\operatorname{dom}(u) \cap \mathrm{E}_{\lambda}^{\circ}\) are stable under \(u\) (prop. 9 of IV, p. 278).

Let \(x \in \text{dom}(u)\) . Let us write \(x = p_{\lambda}(x) + y\) where \(y \in \text{dom}(u) \cap \text{E}_{\lambda}^{\circ}\) . From what precedes, we have

\[
\left\| w (x) - \lambda x \right\| ^ {2} = \left\| p _ {\lambda} (x) \right\| ^ {2} + \left\| w (y) - \lambda y \right\| ^ {2}.\tag{21}
\]

Let V be an open neighborhood of \(\lambda\) not meeting the spectrum of u, and let c > 0 be such that the disk with center \(\lambda\) and of radius c is contained in V. Let \(\mu_{y}\) be the spectral measure of y relative to u. Since y is orthogonal to \(E_{\lambda}\) , the support of \(\mu_{y}\) does not meet V (loc. cit.). We then have

\[
\| w (y) - \lambda y \| ^ {2} = \| u (y) - \lambda y \| ^ {2} = \int_ {\mathbf {C}} | t - \lambda | ^ {2} d \mu_ {y} (t) \geqslant c ^ {2} \| y \| ^ {2}.\tag{22}
\]

We conclude from (21) and (22) that \(\|w(x)-\lambda x\|^{2}\geqslant\inf(c^{2},1)\|x\|^{2}\), since w is self-adjoint and \(\lambda\in R\). The conclusion then follows from Lemma 5 of IV, p. 248.

COROLLARY. --- Let \(u\) be a self-adjoint partial operator on E and \(v\) a compact Hermitian endomorphism of E. We have \(\mathrm{Sp}_{\mathrm{e}}(u + v) = \mathrm{Sp}_{\mathrm{e}}(u)\) .

This follows immediately from the theorem.

